\SourceLeaf{77}
\noindent\hspace*{\parindent}Zweitens hat man aber auch:
\[
\begin{aligned}\Omega &=-\log\Bigl(\dfrac{Y-X}{y-x} . \dfrac{y}{Y}\Bigr)\\ &=-\log\left\{\dfrac{1+\alpha(y+x)+\beta(yy+yx+xx)+\text{etc.}}{1+\alpha y+\beta yy+\gamma y^3+\text{etc.}}\right\}
\end{aligned}
\]
\[
=-\log\left\{1+\dfrac{\alpha x+\beta(yx+xx)+\gamma(yyx+yxx+x^3)+\text{etc.}}{1+\alpha y+\beta yy+\gamma y^3+\text{etc.}}\right\}.
\]

\noindent\hspace*{\parindent}Dieser zweite Ausdruck zeigt, dass $\Omega$ gar keine Potenzen von $y$ mit negativen Exponenten
enthält; es muss also
\[
\dfrac{X}{y}+(2,1)\dfrac{X^2}{2y}+(3,2)\dfrac{X^3}{3y}+\text{etc.}-\dfrac{x}{y}
\]
\noindent identisch $=0$ werden oder
\[
x=X+\tfrac{1}{2}(2,1)X^2+\tfrac{1}{3}(3,2)X^3+\text{etc.},
\]
\noindent welches die Reversionsformel ist.

Eben so folgen hieraus die Coefficienten der Reihen, die $x^2$, $x^3$ etc. oder jede positive ganze
Potenz von $x$ durch $X$ ausdrücken, und durch leichte Kunstgriffe kann man es auch auf gebrochene
Exponenten ausdehnen. Für negative Exponenten müssten noch andere etwas künstlichere Betrachtungen
\mbox{hinzukommen. .\,.\,.\,.\,.\,.\,.}



\Needspace{10\baselineskip}\begin{center}
\textbf{REVERSIO SERIERUM.}
\end{center}

\[
x-a'xx-a''x^3-a'''x^4-a^{\mathrm{IV}}x^5-.\,.\,.\,.=y
\]
\[
\begin{aligned}x^m={}&\sum a^{\prime\lambda'}a^{\prime\prime\lambda''}a^{\prime\prime\prime\lambda'''}.\,.\,. y^{m+\nu}\,\\
&\dfrac{m(\mu+\nu+m-1)(\mu+\nu+m-2)(\mu+\nu+m-3).\,.\,.(\nu+m+1)}
{1.2.3.\,.\,.\lambda'.1.2.3.\,.\,.\lambda''.1.2.3.\,.\,.\lambda'''\,\text{etc.}}\,,
\end{aligned}
\]

\noindent wo
\[
\begin{array}{l}
\mu=\lambda'+\lambda''+\lambda'''+\text{etc.}\\[4pt]
\nu=\lambda'+2\lambda''+3\lambda'''+\text{etc.}
\end{array}
\]

\noindent und $\lambda'$, $\lambda''$, $\lambda'''$ etc. alle möglichen Combinationen von Werthen erhalten,
die weder gebrochen noch negativ sind.

